\section*{3. Salary}
\label{sec:salary}

\subsection*{Dataset structure, provenance, and coverage}

The salary dataset is derived from the deterministic table parser dataset (\texttt{extracted\_data\_salary\_v2.csv}), filtered to high-confidence tiers A and B (\SalTierABShare\ of all extracted rows). Every salary amount is traced directly to a grid cell in the source document. No manual or rule-based QA correction layer applies to salary; data quality is ensured through deterministic parsing, coordinate verification, and confidence filtering.

The tier A+B salary dataset contains \SalTotalRows\ rows, where one row corresponds to a jobgroup -- step -- worker-type -- age -- education combination in a given CAO contract. Table~\ref{tab:salary_structure} summarises the basic structure and full-time hours distribution, while Figure~\ref{fig:salary_rows_by_year} shows how the number of CAOs with salary files evolves over contract start years. Coverage increases quickly from 2008 onwards, peaking at \SalPeakCaoCount\ CAOs contributing salary tables in \SalPeakCaoYear. The drop in 2025--2026 reflects right-censoring.

\begin{table}[H]
    \centering
    \caption{Basic structure of the salary dataset (Parser v2, Confidence Tiers A+B)}
    \label{tab:salary_structure}
    \begin{tabular}{lr}
    \hline\hline
    Characteristic & Value \\
    \hline
    Total salary rows & 343,111 \\
    Unique files (CAO $\times$ PDF) & 2,321 \\
    Unique CAOs & 223 \\
    Distinct job groups & 3,763 \\
    Distinct steps & 10,672 \\
    Distinct worker types & 643 \\
    Distinct age groups & 482 \\
    Distinct education categories & 4 \\
    Share of rows with full-time hours observed & 54\% \\
    Mean full-time weekly hours & 37.63 \\
    Median full-time weekly hours & 38 \\
    \hline\hline
\end{tabular}

\end{table}

\begin{figure}[H]
    \centering
    \includegraphics[width=\textwidth]{figures/salary/salary_boolean_shares_by_contract_year.png}
    \caption{Share of rows with TTW and entry flags, and number of salary rows by contract start year}
    \label{fig:salary_rows_by_year}
\end{figure}

The flag structure in Figure~\ref{fig:salary_rows_by_year} indicates how often rows are marked as interim amendments (\texttt{ttw}) or as entry wages (\texttt{is\_entry}). TTW rows account for between \SalTTWShareLow\ and \SalTTWShareHigh\ of salary rows in mature years, with substantial variation across cohorts, while entry wages make up a consistently small share (\SalEntryShareLow--\SalEntryShareHigh\ of rows per year).

\subsection*{Salary levels and wage growth}

All salary variables in this subsection are normalized to gross monthly EUR. The conversion uses the detected pay unit in the source text: monthly amounts are kept unchanged; four-weekly amounts are multiplied by \(13/12\); hourly wages are multiplied by weekly hours times \(52/12\) (using the CAO's documented full-time hours basis); daily wages by \(5 \times 4.33\); weekly wages by \(52/12\); and annual wages are divided by 12. Salaries are capped at the statutory minimum wage floor and at EUR~50{,}000 per month.

Throughout this subsection, observations are weighted so that each CAO contributes equally within each cohort year, regardless of how many salary scale entries it contains. For contract-year figures, observations receive weights equal to the inverse of the number of observations in their CAO for that contract year. The same CAO-equal weighting logic applies to salary start years and wage increases.

\subsection*{Timeline structure and salary points}

Each salary row can contain multiple timeline points that track how the wage for a given jobgroup--step evolves across a file. Figure~\ref{fig:salary_points_per_row} summarises the amount of salary points at the row level. The median number of positive salary points per row (among rows with at least one) fluctuates between \SalPointsPerRowLow\ and \SalPointsPerRowHigh\ across mature contract cohorts.

\begin{figure}[H]
    \centering
    \includegraphics[width=\textwidth]{figures/salary/salary_points_per_row_by_year.png}
    \caption{Average number of salary timeline points per row by contract start year}
    \label{fig:salary_points_per_row}
\end{figure}

\subsection*{Salary levels}

Figure~\ref{fig:avg_salary_level} plots the CAO-equal-weighted, band-eligible average salary level by salary start year (blue line); in 2025 and 2026 it declines due to right-censoring. As a simpler, unweighted cross-check computed directly on every tier A+B wage-scale row (not CAO-balanced, so it is not directly comparable to the plotted line), the flat average monthly salary rises from \SalLevelSalaryYearEarly\ at the end of the 2000s to \SalLevelSalaryYearLate\ in 2022--2023, reaching \SalLevelSalaryYearRecent\ for the most recent cohorts.

\begin{figure}[H]
    \centering
    \includegraphics[width=\textwidth]{figures/salary/salary_amount_monthly_eur_band_eligible_by_salary_year.png}
    \caption{Whisker plot of the average monthly salary by salary start year (blue line).}
    \label{fig:avg_salary_level}
\end{figure}

Figure~\ref{fig:avg_salary_level_contract_year} plots the same CAO-equal-weighted, band-eligible average by contract start year (blue line); the median is consistently lower but tracks the mean closely, while the overall spread widens over time. The equivalent unweighted flat average across all tier A+B rows (see note above) rises from \SalLevelContractYearEarly\ in \SalLevelContractYearEarlyYear\ to \SalLevelContractYearLate\ in \SalLevelContractYearLateYear\ and \SalLevelContractYearLatest\ in \SalLevelContractYearLatestYear.

\begin{figure}[H]
    \centering
    \includegraphics[width=\textwidth]{figures/salary/salary_amount_monthly_eur_band_eligible_by_contract_year.png}
    \caption{Whisker plot of the average monthly salary by contract start year (blue line).}
    \label{fig:avg_salary_level_contract_year}
\end{figure}

Figure~\ref{fig:avg_salary_level_contract_year_latest} plots the average salary level in the latest CAO view by calendar year (blue line). Compared with the standard contract-year snapshot, the latest CAO view produces a smoother progression in recent cohorts and avoids end-of-sample boundary noise.

\begin{figure}[H]
    \centering
    \includegraphics[width=\textwidth]{figures/salary/salary_amount_monthly_eur_band_eligible_by_contract_year_latest_cao_view.png}
    \caption{Whisker plot of average monthly salary by calendar year in the latest CAO view (blue line).}
    \label{fig:avg_salary_level_contract_year_latest}
\end{figure}

\subsection*{Wage increases}

Figure~\ref{fig:salary_increase_merged_salary_year} plots the average general wage increase by salary year using the merged series (which prefers textual percentage increases when reported, and otherwise computes percentage changes between timeline steps). Average wage increases hover around \SalIncreaseDecadeAvg\ throughout the 2010s, rise after 2021, and peak at \SalIncreasePeakPct\ in \SalIncreasePeakYear\ during the inflationary surge.

\begin{figure}[H]
    \centering
    \includegraphics[width=\textwidth]{figures/salary/salary_increase_merged_pref_csv_by_salary_year.png}
    \caption{Whisker plot of merged average wage increase by salary year (blue line).}
    \label{fig:salary_increase_merged_salary_year}
\end{figure}

Figure~\ref{fig:salary_increase_comparison} compares three alternative wage increase series: calculated differences only (Diff only), textual mentions only (CSV only), and the merged series. The three series track each other closely.

\begin{figure}[H]
    \centering
    \includegraphics[width=\textwidth]{figures/salary/salary_increase_series_comparison_by_year.png}
    \caption{Comparison of wage increase series by salary start year across Diff only, CSV only, and Merged.}
    \label{fig:salary_increase_comparison}
\end{figure}

Figure~\ref{fig:salary_increase_merged_salary_year_latest} presents the merged wage increase measure in the latest CAO view, while Figure~\ref{fig:salary_increase_contract_year} plots the contract start year view. Finally, Figure~\ref{fig:salary_increase_shift_contract_year} displays the contract-to-contract shift in wage increases, capturing the pronounced acceleration during 2022--2023.

\begin{figure}[H]
    \centering
    \includegraphics[width=\textwidth]{figures/salary/salary_increase_merged_pref_csv_by_salary_year_latest_cao_view.png}
    \caption{Merged average wage increase by salary year in latest CAO view (blue line).}
    \label{fig:salary_increase_merged_salary_year_latest}
\end{figure}

\begin{figure}[H]
    \centering
    \includegraphics[width=\textwidth]{figures/salary/salary_increase_percent_by_contract_year.png}
    \caption{Whisker plot of merged wage increase by contract start year (blue line).}
    \label{fig:salary_increase_contract_year}
\end{figure}

\begin{figure}[H]
    \centering
    \includegraphics[width=\textwidth]{figures/salary/salary_increase_shift_by_new_file_year.png}
    \caption{Average shift in wage increases between consecutive contracts by contract start year.}
    \label{fig:salary_increase_shift_contract_year}
\end{figure}

